% !TEX spellcheck = en_US



% ~~~~~~~~~~~~~~~~~~~~~~~~~~~~~~~~~~~~~~~ V
% (NC) 2026 Haluk Bingol
% github.com/halukbingol/LaTeX-Templates
%
% Licensees may copy, distribute, display, and perform the work and make derivative works 
% and remixes based on it only for non-commercial purposes. 
% ~~~~~~~~~~~~~~~~~~~~~~~~~~~~~~~~~~~~~~~ A




% =====================================================================
%  Java on the Command Line -- Sophomore Level: Projects, Libraries and Builds
%  ---------------------------------------------------------------------
%  Built on hbCisLaTeXTemplate.tex with the libraries in ../../hblib/.
%  Programs, terminal sessions and their real outputs are in codeoutput/:
%     codeoutput/<Name>.java  a program        codeoutput/<Name>.sh   a session
%     codeoutput/<Name>/      the session's project folder
%     codeoutput/<Name>.txt   what it printed, made by:  bash run_examples.sh
% =====================================================================

\documentclass[aspectratio=169,10pt,compress]{beamer}

% ~~~~~~~~~~~~~~~~~~~~~~~~~~~~~~~~~~~~~~~~~~~~~~~~~~~~~~~~~~~~~~~~~ V hblib
	\usepackage{../../hblib/hblibPresentation} % lib presentation: theme, i/N, \hSection, algorithm2e, ...
	\usepackage{../../hblib/hblibCodeoutput} % lib code + output: \lstcode, \lstterminal, ...
	\usepackage{../../hblib/hblibDatastructures} % lib data structures: \hString, \hStack, \hQueue
% ~~~~~~~~~~~~~~~~~~~~~~~~~~~~~~~~~~~~~~~~~~~~~~~~~~~~~~~~~~~~~~~~~ A hblib


% xcolor, array (and the L{width} column), listings, tikz and algorithm2e come from the hblib libraries
\usepackage[T1]{fontenc}
\usepackage{lmodern}
\usepackage{booktabs}

\definecolor{gitred}{RGB}{240,80,50}
\newcommand{\win}{\textbf{Windows}}
\newcommand{\mac}{\textbf{macOS}}
\lstdefinelanguage{CSV}{}

% ---------------------------------------------------------------------
\title{Java on the Command Line}
\subtitle{Sophomore Level: Projects, Libraries and Builds}
\author[Bingol]{Haluk O. Bingol}

\institute[FCIS]{
    Faculty of Computer and Information Sciences\\
    Yeditepe University
}
\date{
	{\tiny \hbTimeStamp}
}

\begin{document}




% ~~~~~~~~~~~~~~~~~~~~~~~~~~~~~~~~~~~~~~ V frm
\begin{frame}[t,fragile]
  \titlepage
\end{frame}
% ~~~~~~~~~~~~~~~~~~~~~~~~~~~~~~~~~~~~~~ A frm




% ~~~~~~~~~~~~~~~~~~~~~~~~~~~~~~~~~~~~~~ V frm
\begin{frame}[t,fragile] \frametitle{Outline}
  \begin{columns}[T]
    \begin{column}{0.48\textwidth}
      \tableofcontents[sections={1-3}, hideallsubsections]
    \end{column}
    \begin{column}{0.48\textwidth}
      \tableofcontents[sections={4-6}, hideallsubsections]
    \end{column}
  \end{columns}
\end{frame}
% ~~~~~~~~~~~~~~~~~~~~~~~~~~~~~~~~~~~~~~ A frm

% ~~~~~~~~~~~~~~~~~~~~~~~~~~~~~~~~~~~~~~ V
\hSection[Project]{A Real Project}{}{}
% ~~~~~~~~~~~~~~~~~~~~~~~~~~~~~~~~~~~~~~ A

% ~~~~~~~~~~~~~~~~~~~~~~~~~~~~~~~~~~~~~~ V
\subsection[Layout]{Folder layout}
% ~~~~~~~~~~~~~~~~~~~~~~~~~~~~~~~~~~~~~~ A




% ~~~~~~~~~~~~~~~~~~~~~~~~~~~~~~~~~~~~~~ V frm
\begin{frame}[t,fragile] \frametitle{A library and an application}
  \begin{columns}[T]
    \begin{column}{0.44\textwidth}
      \lstterminal[basicstyle=\ttfamily\tiny]{codeoutput/Layout.txt}
    \end{column}
    \begin{column}{0.54\textwidth}
      \begin{itemize}\small
        \item \texttt{util}: a small \textbf{library}, \texttt{edu.yeditepe.util.Table}, that prints aligned text tables
        \item \texttt{app}: the \textbf{application} \texttt{edu.yeditepe.books}; it reads a book list and prints it with \texttt{Table}
        \item \texttt{app/res}: \textbf{resources}, files the program reads at run time
        \item \texttt{build}: generated \texttt{.class} files; \texttt{dist}: what we ship
        \item \texttt{build.sh} / \texttt{build.bat}: the build as a script
      \end{itemize}
      \begin{block}{Rule}
        \small Keep \textbf{sources} and \textbf{generated} files apart: you can delete \texttt{build} and \texttt{dist} at any time,
        and never put them in Git.
      \end{block}
    \end{column}
  \end{columns}
\end{frame}
% ~~~~~~~~~~~~~~~~~~~~~~~~~~~~~~~~~~~~~~ A frm




% ~~~~~~~~~~~~~~~~~~~~~~~~~~~~~~~~~~~~~~ V frm
\begin{frame}[t,fragile] \frametitle{The library}
  \lstcode[basicstyle=\ttfamily\fontsize{4.6}{5.4}\selectfont]{codeoutput/Project/util/src/edu/yeditepe/util/Table.java}{Java}
\end{frame}
% ~~~~~~~~~~~~~~~~~~~~~~~~~~~~~~~~~~~~~~ A frm




% ~~~~~~~~~~~~~~~~~~~~~~~~~~~~~~~~~~~~~~ V frm
\begin{frame}[t,fragile] \frametitle{The application: \texttt{main} and the data}
  \lstcode[basicstyle=\ttfamily\tiny]{codeoutput/Project/app/src/edu/yeditepe/books/Book.java}{Java}
  \lstcode[basicstyle=\ttfamily\tiny]{codeoutput/Project/app/src/edu/yeditepe/books/Main.java}{Java}
  \lstcode[basicstyle=\ttfamily\tiny, numbers=none]{codeoutput/Project/app/res/books.csv}{CSV}
\end{frame}
% ~~~~~~~~~~~~~~~~~~~~~~~~~~~~~~~~~~~~~~ A frm




% ~~~~~~~~~~~~~~~~~~~~~~~~~~~~~~~~~~~~~~ V frm
\begin{frame}[t,fragile] \frametitle{The application: reading a resource}
  \lstcode[basicstyle=\ttfamily\tiny]{codeoutput/Project/app/src/edu/yeditepe/books/Catalog.java}{Java}

  \vspace{0.1em}
  \small \texttt{getResourceAsStream("/books.csv")} searches the \textbf{class path}, not the current folder,
  so the file is found in a folder \emph{or} inside a JAR.
\end{frame}
% ~~~~~~~~~~~~~~~~~~~~~~~~~~~~~~~~~~~~~~ A frm

% ~~~~~~~~~~~~~~~~~~~~~~~~~~~~~~~~~~~~~~ V
\subsection[By hand]{Building by hand}
% ~~~~~~~~~~~~~~~~~~~~~~~~~~~~~~~~~~~~~~ A




% ~~~~~~~~~~~~~~~~~~~~~~~~~~~~~~~~~~~~~~ V frm
\begin{frame}[t,fragile] \frametitle{Building step by step}
  \lstterminal[basicstyle=\ttfamily\tiny]{codeoutput/ByHand.txt}

  \vspace{0.1em}
  \small
  \texttt{-sourcepath app/src}: \texttt{javac} finds \texttt{Catalog} and \texttt{Book} by itself. \texttt{-cp} names the library
  when compiling \emph{and} when running; \texttt{app/res} is on the class path so \texttt{getResourceAsStream} finds \texttt{books.csv}.
  \win: \texttt{-cp build\textbackslash app;app\textbackslash res;dist\textbackslash lib\textbackslash table.jar}.
\end{frame}
% ~~~~~~~~~~~~~~~~~~~~~~~~~~~~~~~~~~~~~~ A frm

% ~~~~~~~~~~~~~~~~~~~~~~~~~~~~~~~~~~~~~~ V
\hSection[javac]{Compiler Options}{}{}
% ~~~~~~~~~~~~~~~~~~~~~~~~~~~~~~~~~~~~~~ A

% ~~~~~~~~~~~~~~~~~~~~~~~~~~~~~~~~~~~~~~ V
\subsection[Lint]{Warnings}
% ~~~~~~~~~~~~~~~~~~~~~~~~~~~~~~~~~~~~~~ A




% ~~~~~~~~~~~~~~~~~~~~~~~~~~~~~~~~~~~~~~ V frm
\begin{frame}[t,fragile] \frametitle{Warnings: \texttt{-Xlint}}
  \begin{columns}[T]
    \begin{column}{0.46\textwidth}
      \lstcode[basicstyle=\ttfamily\tiny]{codeoutput/Lint/Old.java}{Java}
    \end{column}
    \begin{column}{0.52\textwidth}
      \uncover<2->{\lstterminal[basicstyle=\ttfamily\tiny]{codeoutput/Lint.txt}}
    \end{column}
  \end{columns}

  \vspace{0.1em}
  \small By default \texttt{javac} only prints a short \emph{Note}; \texttt{-Xlint:all} names every problem.
  Add \texttt{-Werror} to treat warnings as errors in a build.
\end{frame}
% ~~~~~~~~~~~~~~~~~~~~~~~~~~~~~~~~~~~~~~ A frm

% ~~~~~~~~~~~~~~~~~~~~~~~~~~~~~~~~~~~~~~ V
\subsection[Encoding]{Source encoding}
% ~~~~~~~~~~~~~~~~~~~~~~~~~~~~~~~~~~~~~~ A




% ~~~~~~~~~~~~~~~~~~~~~~~~~~~~~~~~~~~~~~ V frm
\begin{frame}[t,fragile] \frametitle{Source file encoding: \texttt{-encoding}}
  \lstterminal[basicstyle=\ttfamily\tiny]{codeoutput/Encoding.txt}

  \vspace{0.1em}
  \small
  The file was saved in the old Turkish encoding ISO-8859-9. Since Java 18 \texttt{javac} assumes \textbf{UTF-8} on every system, so
  it rejects the bytes until we name the encoding, or convert the file (\texttt{iconv}; on \win{} use the editor's \emph{Save with encoding}). Older JDKs on Windows assumed Windows-1254 instead: the classic
  ``works on my machine'' bug. \hRemark{Save sources as UTF-8 and pass} \texttt{-encoding UTF-8} \hRemark{in builds.}
\end{frame}
% ~~~~~~~~~~~~~~~~~~~~~~~~~~~~~~~~~~~~~~ A frm

% ~~~~~~~~~~~~~~~~~~~~~~~~~~~~~~~~~~~~~~ V
\subsection[Release]{Target versions}
% ~~~~~~~~~~~~~~~~~~~~~~~~~~~~~~~~~~~~~~ A




% ~~~~~~~~~~~~~~~~~~~~~~~~~~~~~~~~~~~~~~ V frm
\begin{frame}[t,fragile] \frametitle{Compiling for an older Java: \texttt{--release}}
  \begin{columns}[T]
    \begin{column}{0.42\textwidth}
      \lstcode[basicstyle=\ttfamily\tiny]{codeoutput/Release/Modern.java}{Java}
    \end{column}
    \begin{column}{0.56\textwidth}
      \uncover<2->{\lstterminal[basicstyle=\ttfamily\tiny]{codeoutput/Release.txt}}
    \end{column}
  \end{columns}

  \vspace{0.1em}
  \small \texttt{--release N} checks the language \textbf{and} the API of Java N and writes class files that Java N can run
  (version 52 = Java 8, 55 = 11, 61 = 17, 65 = 21, 69 = 25).
  Java 8 is \emph{obsolete} as a target: \texttt{-Xlint:-options} hides that warning.
\end{frame}
% ~~~~~~~~~~~~~~~~~~~~~~~~~~~~~~~~~~~~~~ A frm

% ~~~~~~~~~~~~~~~~~~~~~~~~~~~~~~~~~~~~~~ V
\hSection[Class path]{The Class Path in Depth}{}{}
% ~~~~~~~~~~~~~~~~~~~~~~~~~~~~~~~~~~~~~~ A

% ~~~~~~~~~~~~~~~~~~~~~~~~~~~~~~~~~~~~~~ V
\subsection[Order]{Order, wildcards, CLASSPATH}
% ~~~~~~~~~~~~~~~~~~~~~~~~~~~~~~~~~~~~~~ A




% ~~~~~~~~~~~~~~~~~~~~~~~~~~~~~~~~~~~~~~ V frm
\begin{frame}[t,fragile] \frametitle{Order, wildcards and the \texttt{CLASSPATH} variable}
  \begin{columns}[T]
    \begin{column}{0.34\textwidth}
      \lstcode[basicstyle=\ttfamily\tiny]{codeoutput/ClassPath/App.java}{Java}
    \end{column}
    \begin{column}{0.64\textwidth}
      \uncover<2->{\lstterminal[basicstyle=\ttfamily\tiny]{codeoutput/ClassPath.txt}}
    \end{column}
  \end{columns}

  \vspace{0.1em}
  \small Two JARs with the same class: the \textbf{first} on the class path wins. \texttt{lib/*} means every JAR in \texttt{lib}
  (quote it on macOS; Windows: \texttt{-cp ".;lib\textbackslash *"}). A \texttt{CLASSPATH} variable \textbf{replaces} the default \texttt{.},
  so \texttt{App} itself disappears; prefer \texttt{-cp} and leave \texttt{CLASSPATH} unset.
\end{frame}
% ~~~~~~~~~~~~~~~~~~~~~~~~~~~~~~~~~~~~~~ A frm

% ~~~~~~~~~~~~~~~~~~~~~~~~~~~~~~~~~~~~~~ V
\hSection[Packaging]{Packaging}{}{}
% ~~~~~~~~~~~~~~~~~~~~~~~~~~~~~~~~~~~~~~ A

% ~~~~~~~~~~~~~~~~~~~~~~~~~~~~~~~~~~~~~~ V
\subsection[Class-Path]{The Class-Path attribute}
% ~~~~~~~~~~~~~~~~~~~~~~~~~~~~~~~~~~~~~~ A




% ~~~~~~~~~~~~~~~~~~~~~~~~~~~~~~~~~~~~~~ V frm
\begin{frame}[t,fragile] \frametitle{Executable JAR with a library: \texttt{Class-Path}}
  \lstterminal[basicstyle=\ttfamily\tiny]{codeoutput/Manifest.txt}

  \vspace{0.1em}
  \small
  \texttt{Class-Path: lib/table.jar} is relative to the \textbf{JAR's folder}, not the current folder: ship \texttt{books.jar} with \texttt{lib/};
  alone it fails. \texttt{manifest.txt} must end with a newline, or its last line is silently ignored.
\end{frame}
% ~~~~~~~~~~~~~~~~~~~~~~~~~~~~~~~~~~~~~~ A frm

% ~~~~~~~~~~~~~~~~~~~~~~~~~~~~~~~~~~~~~~ V
\subsection[Fat JAR]{Fat JARs}
% ~~~~~~~~~~~~~~~~~~~~~~~~~~~~~~~~~~~~~~ A




% ~~~~~~~~~~~~~~~~~~~~~~~~~~~~~~~~~~~~~~ V frm
\begin{frame}[t,fragile] \frametitle{A ``fat JAR: everything in one file}
  \lstterminal[basicstyle=\ttfamily\tiny]{codeoutput/FatJar.txt}

  \vspace{0.1em}
  \small
  One file is easy to ship, but you repackage other people's code (check licences), files with the same name collide,
  and signed library JARs break. Build tools automate it (Maven Shade, Gradle Shadow).
\end{frame}
% ~~~~~~~~~~~~~~~~~~~~~~~~~~~~~~~~~~~~~~ A frm

% ~~~~~~~~~~~~~~~~~~~~~~~~~~~~~~~~~~~~~~ V
\subsection[Scripts]{Build scripts}
% ~~~~~~~~~~~~~~~~~~~~~~~~~~~~~~~~~~~~~~ A




% ~~~~~~~~~~~~~~~~~~~~~~~~~~~~~~~~~~~~~~ V frm
\begin{frame}[t,fragile] \frametitle{Build scripts for macOS and Windows}
  \begin{columns}[T]
    \begin{column}{0.49\textwidth}
      \lstcode[basicstyle=\ttfamily\tiny]{codeoutput/Project/build.sh}{bash}
    \end{column}
    \begin{column}{0.49\textwidth}
      \lstcode[basicstyle=\ttfamily\tiny]{codeoutput/Project/build.bat}{command.com}
    \end{column}
  \end{columns}
  \small Run with \hCode{sh build.sh} (macOS Terminal, Git Bash) or \hCode{build.bat} (Windows Command Prompt).
  A script makes the build \textbf{repeatable}; a build tool (next section) adds dependencies, tests and conventions.
\end{frame}
% ~~~~~~~~~~~~~~~~~~~~~~~~~~~~~~~~~~~~~~ A frm

% ~~~~~~~~~~~~~~~~~~~~~~~~~~~~~~~~~~~~~~ V
\hSection[Quality]{Testing and Documentation}{}{}
% ~~~~~~~~~~~~~~~~~~~~~~~~~~~~~~~~~~~~~~ A

% ~~~~~~~~~~~~~~~~~~~~~~~~~~~~~~~~~~~~~~ V
\subsection[Assertions]{Assertions}
% ~~~~~~~~~~~~~~~~~~~~~~~~~~~~~~~~~~~~~~ A




% ~~~~~~~~~~~~~~~~~~~~~~~~~~~~~~~~~~~~~~ V frm
\begin{frame}[t,fragile] \frametitle{Assertions are off by default: \texttt{-ea}}
  \begin{columns}[T]
    \begin{column}{0.56\textwidth}
      \lstcode[basicstyle=\ttfamily\tiny]{codeoutput/Assert/test/TableTest.java}{Java}
    \end{column}
    \begin{column}{0.42\textwidth}
      \uncover<2->{\lstterminal[basicstyle=\ttfamily\tiny]{codeoutput/Assert.txt}}
    \end{column}
  \end{columns}

  \vspace{0.1em}
  \small Without \texttt{-ea} (\texttt{-enableassertions}) every \texttt{assert} is skipped and the broken test ``passes''.
  The second assertion found a real bug: \texttt{render} pads the \textbf{last} column too. Real projects use JUnit.
\end{frame}
% ~~~~~~~~~~~~~~~~~~~~~~~~~~~~~~~~~~~~~~ A frm

% ~~~~~~~~~~~~~~~~~~~~~~~~~~~~~~~~~~~~~~ V
\subsection[javadoc]{API documentation}
% ~~~~~~~~~~~~~~~~~~~~~~~~~~~~~~~~~~~~~~ A




% ~~~~~~~~~~~~~~~~~~~~~~~~~~~~~~~~~~~~~~ V frm
\begin{frame}[t,fragile] \frametitle{Documenting the library}
  \lstterminal[basicstyle=\ttfamily\scriptsize]{codeoutput/Docs.txt}

  \vspace{0.1em}
  \small
  \texttt{-sourcepath util/src edu.yeditepe.util} documents a whole package. The warnings list the missing tags: a method comment should
  describe each parameter (\texttt{@param}) and the result (\texttt{@return}). Ship \texttt{docs} with the library or publish it on a web site.
\end{frame}
% ~~~~~~~~~~~~~~~~~~~~~~~~~~~~~~~~~~~~~~ A frm

% ~~~~~~~~~~~~~~~~~~~~~~~~~~~~~~~~~~~~~~ V
\hSection[Tools]{Build Tools}{}{}
% ~~~~~~~~~~~~~~~~~~~~~~~~~~~~~~~~~~~~~~ A

% ~~~~~~~~~~~~~~~~~~~~~~~~~~~~~~~~~~~~~~ V
\subsection[Maven]{Maven}
% ~~~~~~~~~~~~~~~~~~~~~~~~~~~~~~~~~~~~~~ A




% ~~~~~~~~~~~~~~~~~~~~~~~~~~~~~~~~~~~~~~ V frm
\begin{frame}[t,fragile] \frametitle{Build tools: Maven}
  \begin{columns}[T]
    \begin{column}{0.5\textwidth}


% +++++++++++++++++++++++++++++++++++++++ V lst
%: -Lst.
\begin{lstlisting}[style=codesnippet, language=XML, basicstyle=\ttfamily\tiny]
<project>
  <modelVersion>4.0.0</modelVersion>
  <groupId>edu.yeditepe</groupId>
  <artifactId>books</artifactId>
  <version>1.0</version>
  <properties>
    <maven.compiler.release>21</maven.compiler.release>
    <project.build.sourceEncoding>UTF-8</project.build.sourceEncoding>
  </properties>
  <dependencies>
    <dependency>
      <groupId>edu.yeditepe</groupId>
      <artifactId>table</artifactId>
      <version>1.0</version>
    </dependency>
  </dependencies>
</project>
\end{lstlisting}
% +++++++++++++++++++++++++++++++++++++++ A lst


    \end{column}
    \begin{column}{0.46\textwidth}
      \small \texttt{pom.xml} \emph{describes} the project; Maven knows the steps:
      \vspace{0.2em}

      \footnotesize
      \begin{tabular}{@{}ll@{}}
        \toprule
        \textbf{Command} & \textbf{By hand} \\
        \midrule
        \texttt{mvn compile}  & \texttt{javac -d \dots} \\
        \texttt{mvn test}     & run the tests \\
        \texttt{mvn package}  & \texttt{jar --create \dots} \\
        \texttt{mvn clean}    & \texttt{rm -rf target} \\
        \bottomrule
      \end{tabular}

      \vspace{0.3em}
      \small Sources go in \texttt{src/main/java}, resources in \texttt{src/main/resources}; libraries are downloaded from
      Maven Central. Same commands on Windows and macOS.
      \vspace{0.2em}

      \footnotesize\hRemark{Not run here}: Maven needs the internet, which these examples could not reach.
    \end{column}
  \end{columns}
\end{frame}
% ~~~~~~~~~~~~~~~~~~~~~~~~~~~~~~~~~~~~~~ A frm

% ~~~~~~~~~~~~~~~~~~~~~~~~~~~~~~~~~~~~~~ V
\subsection[Gradle]{Gradle}
% ~~~~~~~~~~~~~~~~~~~~~~~~~~~~~~~~~~~~~~ A




% ~~~~~~~~~~~~~~~~~~~~~~~~~~~~~~~~~~~~~~ V frm
\begin{frame}[t,fragile] \frametitle{Build tools: Gradle}
  \begin{columns}[T]
    \begin{column}{0.5\textwidth}


% +++++++++++++++++++++++++++++++++++++++ V lst
%: -Lst.
\begin{lstlisting}[style=codesnippet, language=Java, basicstyle=\ttfamily\scriptsize]
// build.gradle.kts
plugins {
    application
}

java {
    toolchain { languageVersion = JavaLanguageVersion.of(21) }
}

dependencies {
    implementation(files("lib/table.jar"))
}

application {
    mainClass = "edu.yeditepe.books.Main"
}
\end{lstlisting}
% +++++++++++++++++++++++++++++++++++++++ A lst


    \end{column}
    \begin{column}{0.46\textwidth}
      \footnotesize
      \begin{tabular}{@{}ll@{}}
        \toprule
        \textbf{macOS} & \textbf{Windows} \\
        \midrule
        \texttt{./gradlew build}   & \texttt{gradlew build} \\
        \texttt{./gradlew run}     & \texttt{gradlew run} \\
        \texttt{./gradlew jar}     & \texttt{gradlew jar} \\
        \bottomrule
      \end{tabular}

      \vspace{0.4em}
      \small \texttt{gradlew} (the \emph{wrapper}) downloads the right Gradle version, so everyone builds the same way.
      The \texttt{application} plugin also makes start scripts for Windows and macOS.
      \vspace{0.2em}

      \footnotesize\hRemark{Not run here} (needs the internet). Either tool replaces \texttt{build.sh}; the concepts on these slides stay the same.
    \end{column}
  \end{columns}
\end{frame}
% ~~~~~~~~~~~~~~~~~~~~~~~~~~~~~~~~~~~~~~ A frm

% ~~~~~~~~~~~~~~~~~~~~~~~~~~~~~~~~~~~~~~ V
\hSection[Summary]{Summary}{}{}
% ~~~~~~~~~~~~~~~~~~~~~~~~~~~~~~~~~~~~~~ A

% ~~~~~~~~~~~~~~~~~~~~~~~~~~~~~~~~~~~~~~ V
\subsection[Recipes]{Recipes}
% ~~~~~~~~~~~~~~~~~~~~~~~~~~~~~~~~~~~~~~ A




% ~~~~~~~~~~~~~~~~~~~~~~~~~~~~~~~~~~~~~~ V frm
\begin{frame}[t,fragile] \frametitle{Recipes}
  \footnotesize
  \begin{tabular}{@{}L{0.36\textwidth}L{0.58\textwidth}@{}}
    \toprule
    \textbf{Goal} & \textbf{Command (macOS; Windows: \texttt{;} and \texttt{\textbackslash})} \\
    \midrule
    compile a whole project         & \texttt{javac -d build -sourcepath src -encoding UTF-8 src/a/b/Main.java} \\
    \dots\ against a library        & \texttt{javac -d build -cp lib/x.jar -sourcepath src src/a/b/Main.java} \\
    \dots\ for an older Java        & \texttt{javac --release 17 \dots} \\
    with all warnings               & \texttt{javac -Xlint:all -Werror \dots} \\
    library JAR                     & \texttt{jar --create --file lib/x.jar -C build .} \\
    executable JAR + library        & \texttt{jar --create --file app.jar --manifest m.txt -C build . -C res .} \\
    run with libraries              & \texttt{java -cp "build:res:lib/*" a.b.Main} \\
    run tests with assertions       & \texttt{java -ea -cp build:test Tests} \\
    documentation                   & \texttt{javadoc -d docs -sourcepath src a.b} \\
    \bottomrule
  \end{tabular}
\end{frame}
% ~~~~~~~~~~~~~~~~~~~~~~~~~~~~~~~~~~~~~~ A frm

% ~~~~~~~~~~~~~~~~~~~~~~~~~~~~~~~~~~~~~~ V
\subsection[Exercises]{Exercises}
% ~~~~~~~~~~~~~~~~~~~~~~~~~~~~~~~~~~~~~~ A




% ~~~~~~~~~~~~~~~~~~~~~~~~~~~~~~~~~~~~~~ V frm
\begin{frame}[t,fragile] \frametitle{Exercises}
  \begin{enumerate}\small
    \item Add a column ``Pages'' to \texttt{books.csv} and print it; rebuild with \texttt{build.sh} or \texttt{build.bat}.
    \item Fix the padding bug in \texttt{Table.render} so that \texttt{TableTest} passes with \texttt{-ea}.
    \item Turn \texttt{TableTest} into a script step: the build must stop when a test fails (exit code).
    \item Compile the project with \texttt{--release 17}. What changes? Can you go down to \texttt{--release 11}? (records!)
    \item Ship the application as \texttt{dist/} (JAR + \texttt{lib/}) and as one fat JAR; compare sizes and startup.
    \item Write the \texttt{pom.xml} for the project and build it with Maven on your own computer.
  \end{enumerate}
\end{frame}
% ~~~~~~~~~~~~~~~~~~~~~~~~~~~~~~~~~~~~~~ A frm




% ~~~~~~~~~~~~~~~~~~~~~~~~~~~~~~~~~~~~~~ V frm
\begin{frame}[t,fragile]
	\begin{center}
		\vspace{4cm}
		\hRemark{\Huge Questions?}
	\end{center}
\end{frame}
% ~~~~~~~~~~~~~~~~~~~~~~~~~~~~~~~~~~~~~~ A frm



\end{document}
